\section{Motivation and Challenges}
\label{sec:motivation}

The resulting management problem is to decide which KV to retain and
in which tier. For a multi-agent service, the benefit of these decisions
depends on how retained state supports both normal execution and workflow
completion after a failure. Existing fault-tolerant serving and
workflow-aware KV management provide useful mechanisms, but their
objectives leave this joint decision unresolved.

\subsection{Limitations of Existing Approaches}
\label{sec:motivation-limitations}

\begin{figure*}[t]
  \centering
  \captionsetup[subfigure]{font=normalsize,labelfont=normalfont,
    textfont=normalfont,justification=centering,singlelinecheck=false}
  \begin{subfigure}[t]{2.48in}
    \vspace{0pt}
    \begin{minipage}[c][1.65in][c]{\linewidth}
      \centering
      \includegraphics[width=\linewidth]{figures/workflow.pdf}
    \end{minipage}
    \caption{Four-branch workflow}
    \label{fig:motivation-workflow}
  \end{subfigure}\hfill
  \begin{subfigure}[t]{1.86in}
    \vspace{0pt}
    \begin{minipage}[c][1.65in][c]{\linewidth}
      \centering
      \includegraphics[width=\linewidth]{figures/workflow-gaps.pdf}
    \end{minipage}
    \caption{Workflow-maximum gaps}
    \label{fig:motivation-branch-gap}
  \end{subfigure}\hfill
  \begin{subfigure}[t]{2.52in}
    \vspace{0pt}
    \begin{minipage}[c][1.65in][c]{\linewidth}
      \centering
      \includegraphics[width=\linewidth]{figures/recovery-exposure.pdf}
    \end{minipage}
    \caption{Recovery cost and WCT}
    \label{fig:motivation-recovery-exposure}
  \end{subfigure}
  \caption{Workflow timing and conditional recovery impact.
    (a) Four trace-derived parallel agents A1-A4 with schematic Main agents A0/A5.
    Bars preserve timing; dashed tails assume an all-branch join.
    (b) Empirical CDF of workflow-level gaps over 121 workflows. For each parallel group with temporally inferred fan-in, we subtract the shortest agent execution time from the longest, then retain the largest gap per workflow. 
   (c) Each horizontal bar represents one agent (A1--A4) and shows its
extra delay to reconstruct KV after a failure, measured by full
prefill with Qwen2.5-7B-Instruct-1M. The hatched portion is absorbed
by slack; the solid portion is the calculated WCT increase.}

  \label{fig:motivation}
\end{figure*}

\paragraph{Limitation 1: Fault tolerance without workflow-aware priorities.}
Fault-tolerant serving systems preserve state and coordinate recovery
for interrupted requests. D\'ej\`aVu streams KV checkpoints, and GhostServe
uses host-memory parity to reduce reconstruction
work~\cite{strati2024dejavu,jayakody2026ghostserve}.
LUMEN further coordinates checkpoint placement, recovery routing and
capacity restoration according to worker load~\cite{cao2026lumen}.
These mechanisms address state loss and recovery contention, but do not
use multi-agent dependencies to determine which requests' faster
recovery would advance workflow completion. Request recovery cost and
worker load alone cannot express this benefit.

\label{sec:motivation-workflow}
Figure~\ref{fig:motivation}(\subref{fig:motivation-workflow}) illustrates
why workflow position matters. A2 finishes last, 151.6\,s after the fork,
whereas A1 finishes at 60.4\,s. If the following Main stage requires all
four outputs, A1 can absorb 91.2\,s of additional delay before postponing
the join, whereas any delay on A2 postpones it. This is an example from the Claude code traces~\cite{semianalysis2026cctraces}, which contains the workflows that are very similar to the architecture shown in part a of \ref{fig:agent-workflow-examples}. We call the additional
delay a branch can absorb without postponing the required output its
\emph{slack}. 

This situation where agents with parallel execution have different completion times is very common in multi-agent applications, and the time differences can be significant. We analyze the Claude code trace~\cite{semianalysis2026cctraces} related information and visualize this time difference in Figure~\ref{fig:motivation}(\subref{fig:motivation-branch-gap}).
This figure reports each workflow's maximum longest-minus-shortest branch-duration gap,
with a median of 138\,s and a 90th percentile of 1,173\,s. When the Main agent waits for all required branches, differences in branch completion times can leave shorter branches with slack. Additional delays within this slack do not postpone the join or
increase WCT.

These differences in slack affect how KV reconstruction delays workflow
completion after a failure.
Figure~\ref{fig:motivation}(\subref{fig:motivation-recovery-exposure})
shows reconstruction times measured by full prefill with
Qwen2.5-7B-Instruct-1M using sparse
attention~\cite{qwen2025qwen25oneM,jiang2024minference},
and their calculated impact on WCT.
We consider each agent's interruption separately, with other execution
unchanged and the Main agent waiting for all four branches.
A1's 33.2\,s reconstruction delay is fully absorbed by its 91.2\,s slack,
leaving WCT unchanged. In contrast, A2 has no slack, so its 32.6\,s
reconstruction delay increases WCT by the same amount.
Thus, similar reconstruction costs can have very different effects
on workflow completion. Recovery priorities should therefore account
for a request's position in the workflow, in addition to its
reconstruction cost.

\paragraph{Limitation 2: Workflow-aware KV management ignores failures.}
Workflow-aware serving systems already exploit the dependencies
highlighted above, but to accelerate normal execution. One line of work
hands workflow structure to the scheduler: Parrot exposes inter-request
dependencies for application-level scheduling~\cite{lin2024parrot}, Ayo
orchestrates applications as primitive-level dataflow
graphs~\cite{tan2025ayo}, Autellix and SAGA schedule whole agent programs
rather than individual calls~\cite{luo2025autellix,guo2026saga}, and
Cortex and Murakkab pool and configure resources from the same
structure~\cite{pagonas2025cortex,chaudhry2026murakkab}. Another line
manages KV with it directly: InferCept and Continuum decide whether KV is
kept, swapped or discarded during a tool call, and for how
long~\cite{abhyankar2024infercept,li2025continuum}; TokenCake offloads KV
during tool stalls, prefetches it before the next call, and reserves GPU
memory for critical-path agents~\cite{bian2025tokencake}; KVFlow,
CacheScout and PBKV order eviction and prefetching by predicted upcoming
agent invocations~\cite{pan2025kvflow,zhang2026cachescout,zheng2026pbkv};
Pythia predicts future steps and prompt contents to stage KV ahead of
use~\cite{yu2026pythia}; and TOPAS jointly selects resident prefixes and
admitted requests from remaining-path progress, reuse, and movement or
preemption costs~\cite{ni2026topas}. These policies value retained KV by
the requests expected to read it next, and assume that the worker holding
it keeps running.

Yet failures are not rare, as Figure~\ref{fig:gpu-server-availability}
shows: 584 unavailability events hit 231 of 400 servers in about 349
days. When a server fails, the agent requests running on it are
interrupted together with their GPU-resident KV and any copies swapped
to that host's
memory~\cite{abhyankar2024infercept,pan2025kvflow,bian2025tokencake}.
What remains on other workers or in shared storage are the fixed prompts
and shared prefixes selected by their next expected
read~\cite{pan2025kvflow,yu2026pythia,zhang2026cachescout}; the
interrupted request's history up to the failure is not among them. The
CC/Weka trace bounds how much such surviving prefixes could help: for
follow-up requests, same-agent prefixes cover 98.8\% of main-agent input
and 95.4\% of subagent input, whereas prefixes from other agents cover at
most 4.4\% and 35.7\%, largely overlapping the former
(Figure~\ref{fig:agent-prefix-reuse}). Most of an interrupted request's
input is thus its own agent's history, which resided with that agent on
the failed worker; it must be reconstructed from the prompt and token
log, at the costs of Figures~\ref{fig:motivation-recompute}
and~\ref{fig:motivation-recovery-exposure}, and which agent it lands on
decides whether that time is absorbed by slack or added to WCT. The
root cause is that failure is absent from these systems' model: whether
a KV object is worth keeping, where, and for how long is decided on the
premise that its worker does not fail. Swapping to host memory frees GPU
capacity rather than escaping a failure, shared storage holds completed
contexts for later reuse rather than a running request's history, and a
copy that would be read only after a failure is worth nothing in their
retention scores. So although these systems already hold the information
needed to tell which agent cannot afford this reconstruction---the
dependency graphs, next-call predictions and remaining paths they extract
are exactly what slack is derived from, and TokenCake even uses them to
reserve GPU memory for critical-path agents~\cite{bian2025tokencake}---that
information is never used to decide, before a failure, for which agent,
at which tier, and for how long a resumable copy should be kept.

\begin{figure}[t]
 \centering
 \includegraphics[width=\columnwidth]{figures/agent-prefix-reuse.pdf}
 \caption{Token-weighted potential prefix coverage for 56,405 main-agent
 and 40,332 subagent follow-up requests from CC/Weka. Eligible sources
 are earlier completed requests in the same workflow with the same
 recorded model. Same/other-agent coverage overlaps: other agents add only
 0.01\% and 0.96\% beyond the best same-agent prefix. These are prefix
 opportunities, not runtime cache hits.}
 \label{fig:agent-prefix-reuse}
 \Description{Two groups compare main agents and subagents on follow-up
 requests. Same-agent coverage is 98.8 and 95.4 percent; other-agent
 coverage is 4.4 and 35.7 percent, respectively.}
\end{figure}

\label{sec:motivation-long-context}
The two limitations are mirror images. Fault-tolerant serving decides
\emph{how} to preserve and restore state, but ranks interrupted requests
by recovery cost and worker load. Workflow-aware caching decides
\emph{which} state a workflow will read next, but assumes the worker
holding it survives. Neither answers the question a multi-agent service
faces when a worker fails: which agents' histories must be recoverable
quickly for their workflows to complete on time, and how much of the
finite HBM, DRAM and storage that protection may take from normal reuse.
Combining the two is not a matter of adding slack to a recovery router
or a checkpoint path to a cache. Both objectives draw on the same tiers,
the state each values differs, and both change while agents run.
AegisKV takes this joint decision as its problem. It assumes checkpoint
transport and tiered stores as in the systems above, and contributes the
placement policy---which KV objects each tier holds, on whose behalf,
and when they move---together with the bounded controller and safe
transitions that keep that policy correct online. Realizing it raises
three challenges.

\subsection{Design Challenges}
\label{sec:motivation-challenges}

\paragraph{Challenge 1: Recovery urgency is unknown online and does not add up.}
The examples in Section~\ref{sec:motivation-limitations} obtain slack
from completed traces. Online, the remaining model and tool durations of
every branch are unknown, tool durations are outside the serving
system's control, and the branch that gates the join changes as others
finish. Slack also varies widely: workflow-maximum gaps range from
seconds to over a thousand seconds (median 138\,s, 90th percentile
1,173\,s in Figure~\ref{fig:motivation}(\subref{fig:motivation-branch-gap})),
so no fixed rule separates urgent from tolerant branches, and recency or
reconstruction cost cannot either: A1 and A2 reconstruct in similar time
yet only A2 delays the join. Urgency is not additive. One worker failure
interrupts several co-located agents whose recoveries share readers,
links and target GPUs, and accelerating one of two late branches leaves
their join unchanged. Priorities must therefore be derived from partial
execution, remain useful under prediction error, fall back safely when
estimates are invalid, and rank state without treating request-level
savings as summed WCT reductions.

\paragraph{Challenge 2: Urgency does not say how much of which history belongs in which tier.}
Retaining a history in any tier removes reconstruction: at 1M tokens,
571.21\,s becomes 10.156\,s from shared storage
(Figure~\ref{fig:motivation-recompute}). A faster tier only shortens
retrieval, to 2.663\,s in DRAM and 1.305\,s in HBM, at a very different
price: that history occupies 57.92\,GB, over 70\% of the two A100s'
combined HBM, and every HBM byte spent on checkpoints is taken from
serving. Because interrupted agents' histories are largely unique
(Figure~\ref{fig:agent-prefix-reuse}), sharing does not shrink this
cost, and co-located agents recovering together contend for lower-tier
bandwidth. Whole histories therefore cannot live in the fast tier, but a
history need not move as a unit. Long-context models increasingly select
blocks, tokens or compressed entries per query
(Figure~\ref{fig:sparse-releases}), and their accesses concentrate: in
NOSA-8B, NSA-8B and MiniCPM4.1-8B, about 20\% of historical blocks
receive 87.64\%, 74.76\% and 56.54\% of accesses
(Figure~\ref{fig:design-insights}). A small hot subset in a fast tier
could serve most reads while the complete history stays in cheaper
tiers---but retrospective concentration is not future-query accuracy,
and a missed block must still be retrievable. Placement must therefore
value state per byte and per block, separate retention that avoids
reconstruction from acceleration that shortens retrieval, choose the hot
subset from causal history alone, and credit normal reuse independently
of urgency: high slack must not delete history that later normal calls
would recompute.

\begin{figure}[t]
 \centering
 \includegraphics[width=\columnwidth]{figures/sparse-models.pdf}
 \caption{Selected sparse/hybrid model releases and reported context
 lengths. IDs: 1 MiniCPM4~\cite{minicpm4model};
 2 MiniCPM4.1~\cite{openbmb2025minicpm41};
 3 DeepSeek-V3.2-Exp~\cite{deepseek2025v32exp};
 4 DeepSeek-V3.2~\cite{deepseek2025v32model};
 5 GLM-5~\cite{zai2026glm5model};
 6 MiniCPM-SALA~\cite{openbmb2026sala};
 7 DeepSeek-V4 preview~\cite{deepseek2026v4model};
 8 MiniMax-M3~\cite{minimax2026m3};
 9 GLM-5.2~\cite{zai2026glm52};
 10 LongCat-Flash-Lite-Sparse~\cite{longcat2026sparse}.
 Dates are announcements, or first official repository dates for
 IDs 5, 9 and 10. Markers identify the sparse branch's selection unit,
 including SALA's block branch and V4's compressed sparse branch;
 they do not imply uniform attention across layers. This is a selected
 survey, not an adoption-share or sparsity-rate estimate.}
 \label{fig:sparse-releases}
 \Description{Ten selected sparse or hybrid attention releases between
 June 2025 and July 2026, with reported contexts from 32K to 1M.
 Circles, squares and diamonds distinguish block, token and compressed
 entry selection. This is a literature survey, not a sparsity-rate curve.}
\end{figure}

\begin{figure}[t]
 \centering
 \includegraphics[width=\columnwidth]{figures/kv-access-concentration.pdf}
 \caption{Native KV access concentration for NOSA-8B~\cite{huang2026nosav2},
 NSA-8B (MiniCPM-4-8B-NSA)~\cite{fasa2026minicpm4nsa}, and
 MiniCPM4.1-8B~\cite{openbmb2025minicpm41}. We use the final 512 queries
 of matched 32,704-token prefixes from three PG19 texts. Each query
 counts a historical 64-token block once across heads. Rank the 503
 blocks within each layer/window; average access shares over all
 32 layers, then texts. Bands show text min--max; the diagonal denotes
 uniform access. Compressed/selector state is excluded. This is
 retrospective concentration, not prediction accuracy or recovery speedup.}
 \label{fig:design-insights}
 \Description{Cumulative historical block access curves for NOSA-8B,
 NSA-8B and MiniCPM4.1-8B, averaged over every layer and three texts.
 About 20 percent of blocks account for 87.64, 74.76 and 56.54 percent
 of accesses, respectively; all three curves exceed the uniform baseline.}
\end{figure}

\paragraph{Challenge 3: Placement must follow execution cheaply and move state safely.}
Every generated token extends the unsaved tail, every tool return and
new request changes slack, and every batch changes the HBM that
inference needs. Re-scoring all resident objects on each of these events
is prohibitive at block granularity, yet stale scores keep the wrong
state in the fast tiers. Migration itself consumes the DRAM and storage
bandwidth that recovery and normal reuse depend on, and checkpoint space
in HBM must yield whenever inference needs it. Finally, transfers are
asynchronous and the failure they guard against can strike midway:
releasing a DRAM copy before its storage replacement commits, or
counting an object as local after its owner failed, silently leaves a
checkpoint without a recoverable copy. The controller therefore needs
bounded work per event, a capacity envelope that yields to inference,
and transitions that keep every committed checkpoint with a reachable
owner.

\label{sec:motivation-goals}
These challenges organize the design. Workflow Tracker and the
slack-aware request priority address Challenge~1
(Sections~\ref{sec:workflow-protection} and~\ref{sec:placement-value});
block-level demand, renewable per-agent credit and fixed Peer quotas
address Challenge~2
(Sections~\ref{sec:placement-value}--\ref{sec:sparse-placement}); and the
lazy event loop, AIMD Adapter and versioned transfers address
Challenge~3 (Sections~\ref{sec:placement-plan} and~\ref{sec:adapter}).
